\documentclass[10pt,a4paper]{amsart}
\usepackage[margin=19mm]{geometry}
\usepackage{amsmath,amssymb,mathtools}
\usepackage[scr=rsfs,cal=euler]{mathalfa}
\usepackage{xfrac}
\usepackage[unicode,hidelinks,bookmarks=false]{hyperref}
\newcommand{\ie}{i.e.\ }
\newcommand{\cf}{cf.\ }
\newcommand{\resp}{respectively\ }
\newcommand{\Subsection}{\subsection}
\providecommand{\nobreakdash}{\nobreak\hbox{-}}
\setlength{\emergencystretch}{2em}
\multlinegap0pt
\usepackage{bm}
\usepackage{bbm}
\usepackage{dsfont}
\usepackage{xspace}
\usepackage{cancel}
\usepackage{mathtools}
\usepackage{amsthm}
\makeatletter
\@ifundefined{theorem}{\newtheorem{theorem}{Theorem}}{}
\makeatother
\title[A Functorial Theory of Defects in Abelian Chern-Simons Theor]{A Functorial Theory of Defects in Abelian Chern-Simons Theory}
\author{Daniel Galviz}
\date{Source: arXiv:2609.21414v1 (CC BY 4.0)}
\begin{document}
\maketitle

\noindent\emph{Source and licence.} A Functorial Theory of Defects in Abelian Chern-Simons Theory. Version arXiv:2609.21414v1 (CC BY 4.0), distributed under the Creative Commons Attribution 4.0 International licence, as recorded in the arXiv licence metadata for this exact version and shown on the abstract page https://arxiv.org/abs/2609.21414v1. Full rights notice: licenses/arxiv-2609.21414-CC-BY-4.0.txt.\par\smallskip

\noindent\textit{Editorial scope.} Counted unit: the Theorem whose source label is \texttt{thm:Eq-O} in the article (source0.tex lines 1494--1500, source label \texttt{thm:Eq-O}), delivered together with the article's own complete original proof of it (source0.tex lines 1502--1588, one \texttt{proof} environment of 87 lines). No mathematics is written, repaired or re-derived: statement and proof are the article's own text line for line. Required context retained uncounted: none. Every definition, display and label of those lines is retained unchanged. Omitted from this package and not redistributed: 1--1493, 1501--1501, 1589--5809. Nothing inside a retained excerpt is omitted or altered: every retained body is delivered exactly as the article's lines are written, so no omission receipt of any kind accompanies this package. Citations: the retained excerpts carry no citation command at all, so no bibliography entry of the article is transcribed and no reference is redistributed. Reference adaptation: none was needed and none was made; every reference inside the delivered unit resolves inside it, so no unresolved reference or citation is produced. Printed slips kept literal: no printed word, symbol, hypothesis or number of the counted statement or of its proof was altered. Layout and class: the article's own class is \texttt{amsart} with packages including geometry, fontenc, babel, float, microtype, tikz, tikz-cd, amsmath, amssymb, amsthm, amscd, mathtools, mathrsfs, array; the wrapper is the lane's pinned \texttt{amsart} editorial wrapper at 10pt on A4, which restates the theorem-like environments the retained text needs and adds the article's own macro definitions that the retained text needs (none), with extra wrapper packages (bm, bbm, dsfont, xspace, cancel, mathtools). The delivered numbering is the unit's own. Assets: no retained span contains a figure environment, a graphics-inclusion command, a PDF-page inclusion, a table or a TikZ picture; no image file is copied into this package, the package declares no asset dependency and no image file is redistributed with it. Licence: version arXiv:2609.21414v1 of this article is distributed under the Creative Commons Attribution 4.0 International licence, as recorded in the arXiv licence metadata for this exact version and shown on the abstract page https://arxiv.org/abs/2609.21414v1. A full rights notice is delivered at licenses/arxiv-2609.21414-CC-BY-4.0.txt. Characters: every retained excerpt is pure ASCII, so the delivered engine, XeLaTeX with its default Latin Modern text font, reports no missing character.
\begin{theorem}
\label{thm:Eq-O}
Let $(G,q)$ be a finite quadratic module. Then there is a canonical group isomorphism
$$
\operatorname{Aut}^{\mathrm{br}}(\mathcal C(G,q))\cong \mathcal O(G,q).
$$
\end{theorem}

\begin{proof}
Choose a normalized pointed braided realization $
\mathcal C(G,q)\simeq \operatorname{Vec}^{\,\omega,c}_G,$ where $(\omega,c)$ is an Abelian $3$-cocycle representing the
Eilenberg--Mac Lane class determined by $q$. Thus the simple objects
are $\{X_a\}_{a\in G}$, with
$$
X_a\otimes X_b\simeq X_{a+b},
\qquad
\theta_{X_a}=q(a)\,\operatorname{id}_{X_a}.
$$

Let $F:\mathcal C(G,q)\rightarrow\mathcal C(G,q)$ be a braided ribbon tensor autoequivalence. Since the category is pointed, $F$ permutes the isomorphism classes of simple objects. Hence there is a unique bijection $\phi_F:G\to G$ such that
$$
F(X_a)\cong X_{\phi_F(a)}.
$$
Since $F$ is monoidal,
$$
F(X_{a+b})
\cong
F(X_a\otimes X_b)
\cong
F(X_a)\otimes F(X_b)
\cong
X_{\phi_F(a)+\phi_F(b)},
$$
and therefore $\phi_F(a+b)=\phi_F(a)+\phi_F(b).$ Thus $\phi_F\in\operatorname{Aut}(G)$. Since $F$ is ribbon, it preserves twists, and hence
$$
q(\phi_F(a))=q(a)\qquad \text{for all }a\in G.
$$
Therefore $\phi_F\in\mathcal O(G,q)$. Moreover, braided ribbon
monoidally isomorphic autoequivalences induce the same permutation of
simple objects, so this defines a homomorphism
$$
\Psi:
\operatorname{Aut}^{\mathrm{br}}(\mathcal C(G,q))
\longrightarrow
\mathcal O(G,q),
\qquad
[F]\longmapsto\phi_F.
$$
Conversely, let $\phi\in\mathcal O(G,q)$. Pulling back the chosen Abelian $3$-cocycle along $\phi$ gives $\phi^*(\omega,c).$ Its associated quadratic form is
$$
a\longmapsto c(\phi(a),\phi(a))
=q(\phi(a))
=q(a).
$$
Hence $\phi^*(\omega,c)$ and $(\omega,c)$ determine the same Eilenberg--Mac Lane quadratic class. By the classification of pointed braided categories by quadratic forms, they are Abelian-cohomologous. Consequently there exists a normalized $2$-cochain
$$
\eta_\phi:G\times G\longrightarrow\mathbb C^\times
$$
whose Abelian coboundary identifies $\phi^*(\omega,c)$ with $(\omega,c)$. Using $\eta_\phi$ as tensorator makes the relabelling
$$
X_a\longmapsto X_{\phi(a)}
$$
into a braided ribbon tensor autoequivalence, which we denote by $F_\phi$.

Although the cochain $\eta_\phi$ need not be unique, the resulting braided monoidal isomorphism class of $F_\phi$ is unique. Indeed, if $\eta_\phi$ and $\eta'_\phi$ are two choices, then their ratio
$$
\xi(a,b):=\eta'_\phi(a,b)\eta_\phi(a,b)^{-1}
$$
is a normalized $2$-cocycle. Since both tensorators satisfy the braided compatibility condition, one has
$$
\xi(a,b)=\xi(b,a).
$$
Thus $\xi$ is a symmetric normalized $2$-cocycle. For a finite Abelian group $G$ with coefficients in the divisible group $\mathbb C^\times$, every symmetric normalized $2$-cocycle is a coboundary. Hence the two choices of tensorator give braided monoidally isomorphic functors.

It follows that
$$
\Phi:
\mathcal O(G,q)
\longrightarrow
\operatorname{Aut}^{\mathrm{br}}(\mathcal C(G,q)),
\qquad
\phi\longmapsto[F_\phi],
$$
is well defined. It is a homomorphism: the composite $F_\phi\circ F_\psi$ induces the relabelling $\phi\circ\psi$, and by the uniqueness just proved its braided monoidal isomorphism class is $[F_{\phi\circ\psi}]$.

By construction,
$$
\Psi(\Phi(\phi))=\phi.
$$
Conversely, if $F$ is any braided ribbon tensor autoequivalence and $\phi=\phi_F$, then $F$ and $F_\phi$ induce the same relabelling of simple objects. The same uniqueness argument shows that they are braided ribbon monoidally isomorphic. Therefore
$$
\Phi(\Psi([F]))=[F].
$$
Thus $\Phi$ and $\Psi$ are mutually inverse group isomorphisms.
\end{proof}

\end{document}
